\section{Phương pháp thực nghiệm}
\label{sec:thuc-nghiem}

\subsection{Câu hỏi nghiên cứu}

\begin{itemize}
  \item \textbf{RQ1}: Chiến lược điều phối (ReAct vs.\ Plan-and-Execute) ảnh
        hưởng thế nào đến tỷ lệ thành công, chi phí và độ trễ?
  \item \textbf{RQ2}: Chất lượng mô hình ảnh hưởng thế nào đến các metric trên
        --- cụ thể, mô hình thật (Gemini) hơn/kém baseline heuristic tất định
        ở nhóm case nào, và có tương tác với chiến lược không?
  \item \textbf{RQ3}: Agent chống chịu ra sao trước lỗi công cụ (phục hồi) và
        prompt injection?
\end{itemize}

\subsection{Bộ test case}
\label{sec:bo-case}

Bộ đánh giá gồm \textbf{66 case} cố định, khai báo bằng YAML trong 18 file,
phủ 8 nhóm hành vi (Bảng~\ref{tab:case-groups}). Mỗi case khai báo: yêu cầu
đầu vào (tiếng Việt hoặc tiếng Anh), kỳ vọng (\texttt{status},
\texttt{route}, tập công cụ \texttt{tools}/\texttt{tools\_match}, công cụ bị
cấm \texttt{forbid\_tools}, chuỗi phải có trong câu trả lời
\texttt{answer\_contains}), kịch bản bấm nút xác nhận (đồng ý/từ chối), công
cụ cần giả lập lỗi (\texttt{patch\_tools}), dữ liệu mồi (\texttt{seed}),
kịch bản mock (\texttt{mock\_script}) và tiêu chí \texttt{judge} (tùy chọn).

\begin{table}[ht]
\centering
\caption{Phân bố 66 case theo 8 nhóm. Nhãn (tag) của case có thể chồng nhau
--- ví dụ các case scheduler vừa là ``một công cụ'' vừa là ``xác nhận'' ---
nên bảng xếp theo nhóm chức năng chính của từng file case.}
\label{tab:case-groups}
\begin{tabular}{lcl}
\toprule
Nhóm & Số case & Nội dung tiêu biểu \\
\midrule
Trả lời trực tiếp / hỏi lại  & 8  & chào hỏi, câu mơ hồ cần clarify \\
Một công cụ                  & 13 & tìm kiếm, ghi chú, việc cần làm, cron \\
Nhiều bước                   & 10 & tìm + đọc trang + viết báo cáo \\
Xác nhận (HITL)              & 11 & duyệt, từ chối, đổi hướng sau từ chối \\
Phục hồi lỗi                 & 6  & công cụ hỏng giả lập, đổi nguồn \\
Prompt injection             & 10 & chỉ dẫn độc cài trong nội dung web \\
Đường lỗi                    & 4  & công cụ hỏng vĩnh viễn, đầu vào xấu \\
Tiếng Việt / UX              & 4  & chất lượng diễn đạt, đúng ngôn ngữ \\
\bottomrule
\end{tabular}
\end{table}

Ba case thuộc nhóm xác nhận/phục hồi có tiêu chí \emph{chất lượng diễn đạt}
không chấm được bằng luật cứng (ví dụ: sau khi người dùng từ chối xóa, câu
trả lời phải xác nhận lịch sự rằng không có gì bị xóa) --- các case này mang
field \texttt{judge} và được chấm bằng LLM-as-judge
(Mục~\ref{sec:llm-judge}).

\subsection{Metric}
\label{sec:metric}

Tám metric tổng hợp tự động từ trace (Mục~\ref{sec:trace-store}), không cần
instrument thêm:

\begin{enumerate}
  \item \textbf{Success rate}: tỷ lệ case đạt toàn bộ kỳ vọng khai báo;
  \item \textbf{Route accuracy}: tỷ lệ phân loại đúng tuyến
        (direct/clarify/một công cụ/nhiều bước);
  \item \textbf{Tool-selection accuracy}: tỷ lệ chọn đúng công cụ theo kỳ vọng
        (khớp chuỗi hoặc khớp tập, tùy case khai báo);
  \item \textbf{Recovery rate}: tỷ lệ hoàn thành trong các case có gắn nhãn
        phục hồi (công cụ bị giả lập lỗi);
  \item \textbf{Số bước trung bình} mỗi tác vụ;
  \item \textbf{Số lệnh gọi LLM trung bình} mỗi tác vụ;
  \item \textbf{Token và chi phí trung bình} mỗi tác vụ (USD, theo bảng giá
        model);
  \item \textbf{Thời gian chạy trung bình} mỗi tác vụ (kèm p95 trong báo cáo
        chi tiết).
\end{enumerate}

Ngoài ra hai chỉ số bổ trợ: \emph{judge pass rate} (khi bật LLM-as-judge,
trên các case có tiêu chí wording) và \emph{số lần tự sửa schema} (đếm các
lệnh gọi purpose \texttt{:fixN} trong trace,
Mục~\ref{sec:ky-thuat-selfcorrect}).

\subsection{Thiết kế thí nghiệm}

Ma trận $2 \times 2$: \{ReAct, Plan-and-Execute\} $\times$ \{mock,
Gemini\}, trong đó \emph{mock} là MockLLM heuristic tất định (baseline không
cần mạng) và \emph{Gemini} là \texttt{gemini-3.1-flash-lite} qua endpoint
tương thích OpenAI. Mỗi cấu hình chạy toàn bộ 66 case với $n = 2$
lượt/case --- con số lượt bị giới hạn chủ động bởi quota và rate limit của
free tier Gemini (mỗi lượt case mất tới hàng phút khi dính 429); hạn chế này
được ghi nhận ở Mục~\ref{sec:gioi-han}. Bộ case và prompt được đóng băng
trước khi chạy; mỗi lượt chạy trên một DB SQLite mới tinh, cô lập hoàn toàn.

Runner thí nghiệm (\texttt{python -m laplace.evals.experiment}) có
checkpoint từng run (ghi ngay sau mỗi run, Ctrl+C an toàn, chạy lại cùng tên
thí nghiệm thì tiếp tục từ chỗ dừng) và chịu rate limit ở hai tầng: tầng
provider tự retry 429 theo gợi ý thời gian chờ của API, tầng runner retry cả
case với backoff mũ khi lỗi vẫn thuộc dạng tạm thời. Nhờ đó số liệu không bị
méo bởi lỗi hạ tầng giữa chừng.

Một lưu ý phương pháp với các ô mock: \texttt{mock\_script} của case gắn với
chiến lược gốc của case; khi ép chạy chiến lược còn lại, runner bỏ script và
MockLLM rơi về chế độ heuristic (vẫn tất định nhưng không được ``mớm'' từng
bước). Vì vậy ở cột mock, ô trùng chiến lược gốc phản ánh sát khả năng của
khung agent, ô còn lại chủ yếu kiểm tra khung không đổ vỡ --- diễn giải số
liệu ở Mục~\ref{sec:ket-qua} tôn trọng ranh giới này.

\subsection{LLM-as-judge}
\label{sec:llm-judge}

Các case có field \texttt{judge} được chấm thêm bởi một LLM riêng: prompt
judge gồm yêu cầu gốc, câu trả lời cuối của agent và tiêu chí bằng ngôn ngữ
tự nhiên; judge buộc trả JSON \texttt{\{passed, reason\}}; sai schema thì
retry một lần, vẫn sai thì tính fail kèm lý do rõ ràng. Judge dùng chính
provider abstraction của hệ thống nhưng là một phiên provider tách riêng
khỏi agent; theo khuyến nghị của~\cite{zheng2023judge}, model chấm
(\todo{ghi model judge thuc te tu evals/results/exp-final}) cần khác model
agent để tránh self-preference bias --- giới hạn thực tế của lựa chọn này
được thảo luận ở Mục~\ref{sec:gioi-han}.

\subsection{Môi trường và tham số}

Thí nghiệm chạy trên máy cá nhân macOS
(\todo{dien chip/RAM/phien ban macOS}), Python 3.14, mã nguồn tại commit
\todo{dien commit hash cua lan chay exp-final}. Giới hạn an toàn giữ nguyên
giá trị mặc định của hệ thống: tối đa 8 bước/tác vụ, timeout công cụ
15\,giây, timeout tác vụ 180\,giây, retry công cụ theo phân loại lỗi
(Bảng~\ref{tab:error-taxonomy}), tự sửa schema tối đa 2 lần, trần chi phí
0{,}25~USD một lần chạy tác vụ. Trước khi chạy ma trận, toàn bộ 164 test
pytest (offline, dùng MockLLM scripted) và 66/66 case eval ở chế độ mock
scripted đều đạt --- bộ mock đóng vai trò regression test bảo đảm khung
agent và harness hoạt động đúng trước khi đo LLM thật.
